\chapter{Deep Listening\\การฟังที่ทำให้ความคิดปรากฏ}

การฟังเชิงลึกไม่ใช่ความเงียบระหว่างรอพูด แต่เป็นการให้ความสนใจจน Mentee สามารถได้ยินความคิดของตนเองชัดขึ้น Mentor ที่ฟังดีไม่รีบแปลทุกปัญหาเป็นประสบการณ์ของตน และไม่รีบแก้ความไม่สบายใจด้วยคำแนะนำ

\learningobjectives{
  \item แยกระดับการฟังจากเนื้อหา อารมณ์ และสมมติฐาน
  \item ใช้การสะท้อน ความเงียบ และการสรุปอย่างเหมาะสม
  \item สังเกตสิ่งรบกวนและอคติของตนระหว่างสนทนา
}

\section{ฟังสามชั้น}
ชั้นแรกคือ \textbf{เนื้อหา}: เกิดอะไรขึ้น ชั้นที่สองคือ \textbf{ความหมายและอารมณ์}: สิ่งนี้สำคัญอย่างไรต่อ Mentee ชั้นที่สามคือ \textbf{รูปแบบและสมมติฐาน}: เรื่องใดเกิดซ้ำ สิ่งใดถูกละไว้ และกรอบคิดใดกำหนดทางเลือก การฟังทั้งสามชั้นทำให้ Mentor ถามได้ลึกโดยไม่คาดเดาเกินข้อมูล

\section{พฤติกรรมของการฟัง}
จัดสภาพแวดล้อมให้ปราศจากสิ่งรบกวน ใช้สายตาและท่าทางที่เป็นธรรมชาติ เว้นความเงียบหลังคำตอบ สะท้อนคำสำคัญ และสรุปเพื่อตรวจความเข้าใจ เช่น \enquote{ผมได้ยินว่าท่านไม่ได้กังวลเรื่องเทคนิคเท่ากับกังวลว่าจะชวนทีมเปลี่ยนวิธีทำงาน ผมเข้าใจตรงไหม}

\begin{examplebox}
\textbf{Mentee:} \enquote{โครงการคงไปต่อไม่ได้ เพราะหน่วยงานไม่ตอบ}\\
\textbf{คำตอบที่รีบแก้:} \enquote{ผมรู้จักผู้อำนวยการ เดี๋ยวโทรให้}\\
\textbf{คำตอบที่ฟังก่อน:} \enquote{เมื่อบอกว่าไปต่อไม่ได้ ส่วนใดติดอยู่ และการไม่ตอบทำให้ท่านตีความว่าอย่างไร}\\
คำตอบหลังยังไม่ปฏิเสธการช่วยเชื่อม แต่เปิดพื้นที่ให้เข้าใจปัญหาจริงก่อน
\end{examplebox}

\section{ฟังตนเองไปพร้อมกัน}
Mentor ควรสังเกตแรงกระตุ้นภายใน เช่น อยากแสดงความรู้ อยากให้ Mentee เลือกทางเดียวกับตน หรือรู้สึกไม่อดทนต่อความลังเล การรับรู้นี้ไม่ได้ทำให้อคติหายไป แต่ช่วยไม่ให้อคติกำกับบทสนทนาโดยไม่รู้ตัว

\begin{mentornote}
ก่อนให้คำแนะนำ ลองสรุปสิ่งที่ได้ยินและถามว่า \enquote{มีส่วนสำคัญใดที่ผมยังไม่เข้าใจ} หาก Mentee แก้ความเข้าใจของท่าน แสดงว่าการฟังกำลังทำงาน
\end{mentornote}

\begin{keytakeaway}
Deep Listening สร้างทั้งข้อมูลและความไว้วางใจ เมื่อ Mentee รู้สึกว่าได้รับการฟัง เขาจะกล้าพูดถึงปัญหาจริง ซึ่งเป็นจุดเริ่มของการเปลี่ยนแปลงที่แท้จริง
\end{keytakeaway}

\begin{reflection}
ในการสนทนาครั้งถัดไป บันทึกหลังจบว่า Mentee พูดเรื่องใด อารมณ์หรือคุณค่าใดปรากฏ และสมมติฐานใดควรถามต่อ ระบุหนึ่งช่วงที่ท่านรีบคิดคำตอบแทนการฟัง
\blankline\blankline
\end{reflection}

\chaptersummary{การฟังเชิงลึกผสานความสนใจต่อเนื้อหา ความหมาย และรูปแบบ Mentor ใช้ความเงียบ การสะท้อน และการตรวจความเข้าใจเพื่อช่วยให้ปัญหาจริงปรากฏ โดยตระหนักถึงอคติของตนตลอดเวลา}
